\chapter{Analisi dei Requisiti}

\section{Obiettivi dello stage}

\subsection*{Notazione}
Si farà riferimento ai requisiti secondo leseguenti notazioni:
\begin{itemize}
    \item \textbf{O} per i requisiti obbligatori, vincolanti in quanto obiettivo primario richiesto dal committente;
    \item \textbf{D} per i requisiti desiderabili, non cinvolanti o strettamente necessari, ma dal riconoscibile valore aggiunto;
    \item \textbf{F} per i requisiti facoltativi, rappresenanti valore aggiunto non strettamente competitivo.
\end{itemize}
Le sigle precedentemente indicate saranno seguite da una coppia sequenziale di numeri, identificativo del requisito.

\subsection*{Obiettivi fissati}
Si prevede lo svolgimento dei seguenti obiettivi:
\begin{itemize}
    \item Obbligatori
    \begin{itemize}
        \item O01: Implementazione flusso mediante \nameref{sec:Spring};
        \item O02: Relazione tecnica che descriva le scelte tecniche prese durante il percorso e le motivazioni a riguardo;
        \item \dots
    \end{itemize}
    \item Desiderabili
    \begin{itemize}
        \item D01: Implementazione flusso mediante Apache Flink;
        \item D02: Affrontare un'analisi comparativa usando metriche ben precise e coerenti con il contesto applicativo
    \end{itemize}
    \item Facoltativi
    \begin{itemize}
        \item F01: Effettuare dei test di performance mediante strumenti di misurazione dedicati;
        \item F02: Implementare il flusso anche su tecnologie alternative (Apache Camel, Apache NiFi)
    \end{itemize}
\end{itemize}

\section{Tracciamento dei requisiti}

\section{Tabella dei requisiti}

\begin{table}[H]
\centering
\begin{tabular}{|l|l|l|}
\hline
Requisito & Descrizione & Stato \\ \hline
asd & asd & asd \\ \hline
asd & asd &  \\ \hline
asd &  &  \\ \hline
\end{tabular}
\caption{Tabella requisiti}
\label{}
\end{table}

asdasd


\newpage
